% !TEX program = tectonic
% =============================================================================
%  AULA 4 — Redes sem Fio: Protocolos de acesso ao meio (MAC)
%  Curso: Redes sem Fio  |  Profa. Anelise Munaretto
%  TDMA/FDMA, ALOHA, CSMA/CD, CSMA/CA, RTS/CTS, 802.11e.
%  Compilar: tectonic aula4.tex
% =============================================================================
\documentclass[aspectratio=169,11pt]{beamer}
% =============================================================================
%  PREÁMBULO COMÚN para as aulas de Redes sem Fio (Beamer + TikZ)
%  Incluido com % =============================================================================
%  PREÁMBULO COMÚN para as aulas de Redes sem Fio (Beamer + TikZ)
%  Incluido com % =============================================================================
%  PREÁMBULO COMÚN para as aulas de Redes sem Fio (Beamer + TikZ)
%  Incluido com \input{preamble.tex} en cada aula.
%  Paleta de cores e estilos são definidos aqui uma única vez.
% =============================================================================

\usepackage[T1]{fontenc}
\usepackage{amsmath,amssymb,mathtools}
\usepackage{graphicx}
\usepackage{tikz}
\usetikzlibrary{positioning,arrows.meta,calc,backgrounds,decorations.pathmorphing,fit,shadows,shapes.symbols,shapes.geometric}
\usepackage{pgfplots}
\pgfplotsset{compat=1.16}
\usepackage{tabularx,booktabs,multirow,array}
\usepackage[normalem]{ulem}
\usepackage{hyperref}

% ---------- Paleta de cores própria ----------
\definecolor{aneliseAzul}{RGB}{20,60,110}
\definecolor{aneliseVerde}{RGB}{30,140,70}
\definecolor{aneliseLaranja}{RGB}{215,120,20}
\definecolor{aneliseGris}{RGB}{90,90,90}

% ---------- Tema ----------
\usetheme{Madrid}
\usecolortheme{whale}
\setbeamercolor{structure}{fg=aneliseAzul}
\setbeamercolor{title}{fg=aneliseAzul}
\setbeamercolor{frametitle}{fg=aneliseAzul}
\setbeamercolor{block title}{fg=aneliseAzul,bg=aneliseGris!12!white}
\setbeamercolor{block body}{bg=aneliseGris!7!white}
\hypersetup{colorlinks=true,linkcolor=aneliseAzul,urlcolor=aneliseVerde}

% ---------- Comandos auxiliares ----------
\newcommand{\kurose}{Kurose \& Ross, \emph{Computer Networking: a Top-Down Approach}}
\newcommand{\forouzan}{Forouzan, \emph{Data Communications and Networking}}
\newcommand{\stallings}{Stallings, \emph{Wireless Communications \& Networks}}
\newcommand{\tanenbaum}{Tanenbaum, \emph{Computer Networks}}
\newcommand{\rfc}[1]{\href{https://www.rfc-editor.org/rfc/rfc#1.txt}{RFC~#1}}
% -----------------------------------------------------------------------------
%  Estilos TikZ comuns para desenhar redes (posicionamento relativo)
% -----------------------------------------------------------------------------
\tikzset{
  host/.style={draw,rounded corners=2mm,fill=aneliseAzul!15,inner sep=1.5mm,font=\scriptsize},
  rt/.style={draw,rounded corners=1mm,fill=aneliseLaranja!25,inner sep=1mm,font=\scriptsize},
  nube/.style={draw,cloud,aspect=2.2,fill=aneliseGris!15,inner sep=1.5mm,font=\scriptsize},
  el/.style={draw,rounded corners=1.2mm,fill=aneliseGris!18,inner sep=0.8mm,font=\scriptsize},
  enl/.style={thick},
  enlA/.style={thick,-Stealth},
} en cada aula.
%  Paleta de cores e estilos são definidos aqui uma única vez.
% =============================================================================

\usepackage[T1]{fontenc}
\usepackage{amsmath,amssymb,mathtools}
\usepackage{graphicx}
\usepackage{tikz}
\usetikzlibrary{positioning,arrows.meta,calc,backgrounds,decorations.pathmorphing,fit,shadows,shapes.symbols,shapes.geometric}
\usepackage{pgfplots}
\pgfplotsset{compat=1.16}
\usepackage{tabularx,booktabs,multirow,array}
\usepackage[normalem]{ulem}
\usepackage{hyperref}

% ---------- Paleta de cores própria ----------
\definecolor{aneliseAzul}{RGB}{20,60,110}
\definecolor{aneliseVerde}{RGB}{30,140,70}
\definecolor{aneliseLaranja}{RGB}{215,120,20}
\definecolor{aneliseGris}{RGB}{90,90,90}

% ---------- Tema ----------
\usetheme{Madrid}
\usecolortheme{whale}
\setbeamercolor{structure}{fg=aneliseAzul}
\setbeamercolor{title}{fg=aneliseAzul}
\setbeamercolor{frametitle}{fg=aneliseAzul}
\setbeamercolor{block title}{fg=aneliseAzul,bg=aneliseGris!12!white}
\setbeamercolor{block body}{bg=aneliseGris!7!white}
\hypersetup{colorlinks=true,linkcolor=aneliseAzul,urlcolor=aneliseVerde}

% ---------- Comandos auxiliares ----------
\newcommand{\kurose}{Kurose \& Ross, \emph{Computer Networking: a Top-Down Approach}}
\newcommand{\forouzan}{Forouzan, \emph{Data Communications and Networking}}
\newcommand{\stallings}{Stallings, \emph{Wireless Communications \& Networks}}
\newcommand{\tanenbaum}{Tanenbaum, \emph{Computer Networks}}
\newcommand{\rfc}[1]{\href{https://www.rfc-editor.org/rfc/rfc#1.txt}{RFC~#1}}
% -----------------------------------------------------------------------------
%  Estilos TikZ comuns para desenhar redes (posicionamento relativo)
% -----------------------------------------------------------------------------
\tikzset{
  host/.style={draw,rounded corners=2mm,fill=aneliseAzul!15,inner sep=1.5mm,font=\scriptsize},
  rt/.style={draw,rounded corners=1mm,fill=aneliseLaranja!25,inner sep=1mm,font=\scriptsize},
  nube/.style={draw,cloud,aspect=2.2,fill=aneliseGris!15,inner sep=1.5mm,font=\scriptsize},
  el/.style={draw,rounded corners=1.2mm,fill=aneliseGris!18,inner sep=0.8mm,font=\scriptsize},
  enl/.style={thick},
  enlA/.style={thick,-Stealth},
} en cada aula.
%  Paleta de cores e estilos são definidos aqui uma única vez.
% =============================================================================

\usepackage[T1]{fontenc}
\usepackage{amsmath,amssymb,mathtools}
\usepackage{graphicx}
\usepackage{tikz}
\usetikzlibrary{positioning,arrows.meta,calc,backgrounds,decorations.pathmorphing,fit,shadows,shapes.symbols,shapes.geometric}
\usepackage{pgfplots}
\pgfplotsset{compat=1.16}
\usepackage{tabularx,booktabs,multirow,array}
\usepackage[normalem]{ulem}
\usepackage{hyperref}

% ---------- Paleta de cores própria ----------
\definecolor{aneliseAzul}{RGB}{20,60,110}
\definecolor{aneliseVerde}{RGB}{30,140,70}
\definecolor{aneliseLaranja}{RGB}{215,120,20}
\definecolor{aneliseGris}{RGB}{90,90,90}

% ---------- Tema ----------
\usetheme{Madrid}
\usecolortheme{whale}
\setbeamercolor{structure}{fg=aneliseAzul}
\setbeamercolor{title}{fg=aneliseAzul}
\setbeamercolor{frametitle}{fg=aneliseAzul}
\setbeamercolor{block title}{fg=aneliseAzul,bg=aneliseGris!12!white}
\setbeamercolor{block body}{bg=aneliseGris!7!white}
\hypersetup{colorlinks=true,linkcolor=aneliseAzul,urlcolor=aneliseVerde}

% ---------- Comandos auxiliares ----------
\newcommand{\kurose}{Kurose \& Ross, \emph{Computer Networking: a Top-Down Approach}}
\newcommand{\forouzan}{Forouzan, \emph{Data Communications and Networking}}
\newcommand{\stallings}{Stallings, \emph{Wireless Communications \& Networks}}
\newcommand{\tanenbaum}{Tanenbaum, \emph{Computer Networks}}
\newcommand{\rfc}[1]{\href{https://www.rfc-editor.org/rfc/rfc#1.txt}{RFC~#1}}
% -----------------------------------------------------------------------------
%  Estilos TikZ comuns para desenhar redes (posicionamento relativo)
% -----------------------------------------------------------------------------
\tikzset{
  host/.style={draw,rounded corners=2mm,fill=aneliseAzul!15,inner sep=1.5mm,font=\scriptsize},
  rt/.style={draw,rounded corners=1mm,fill=aneliseLaranja!25,inner sep=1mm,font=\scriptsize},
  nube/.style={draw,cloud,aspect=2.2,fill=aneliseGris!15,inner sep=1.5mm,font=\scriptsize},
  el/.style={draw,rounded corners=1.2mm,fill=aneliseGris!18,inner sep=0.8mm,font=\scriptsize},
  enl/.style={thick},
  enlA/.style={thick,-Stealth},
}

\title[Redes sem Fio]{Redes sem Fio\\ Protocolos de acesso ao meio}
\subtitle{Aula 4}
\author[Anelise Munaretto]{Profa. Anelise Munaretto}
\institute[Curso]{Redes sem Fio --- Ciência da Computação}
\date{2026}

\begin{document}

\begin{frame}[plain]
  \titlepage
\end{frame}

\begin{frame}{Objetivos de aprendizagem}
  Ao final desta aula, você será capaz de:
  \begin{itemize}
    \item \textbf{calcular} a eficiência do Slotted ALOHA ($S=N\,p\,(1-p)^{N-1}$) e do ALOHA puro;
    \item \textbf{explicar} por que o CSMA/CD não funciona em enlaces de rádio;
    \item \textbf{descrever} o problema do terminal oculto e o papel do RTS/CTS e do NAV;
    \item \textbf{comparar} as famílias de protocolos MAC: particionamento, acesso aleatório e passagem de turno;
    \item \textbf{relacionar} os mecanismos do 802.11 (DIFS/SIFS, backoff, PCF) com a operação do CSMA/CA.
  \end{itemize}
  \begin{alertblock}{Pré-requisitos}
    Aula 1 (camadas e arquitetura da Internet) e Aula 2 (características de enlaces sem fio).
  \end{alertblock}
\end{frame}

\section{Enlaces e protocolos de acesso múltiplo}

\begin{frame}{Enlaces de acesso múltiplo e protocolos}
  \textbf{Dois tipos de enlaces:}
  \begin{itemize}
    \item \textbf{Ponto a ponto} (ex.: protocolo PPP)
    \item \textbf{Broadcast} (fio ou meio compartilhado): Ethernet tradicional, 802.11 LAN sem fio
  \end{itemize}
  \begin{block}{Problema central}
    Canal de comunicação único e compartilhado:
    \begin{itemize}
      \item duas ou mais transmissões simultâneas $\rightarrow$ \textbf{interferência}
      \item \textbf{colisão} se um nó recebe duas ou mais sinais ao mesmo tempo
      \item protocolo de acesso múltiplo $=$ algoritmo distribuído que determina quando cada estação pode transmitir
      \item \alert{não há canal out-of-band para coordenação}
    \end{itemize}
  \end{block}
\end{frame}

\begin{frame}{Protocolo ideal de acesso múltiplo}
  Canal de broadcast de taxa $R$ bps. Objetivos:
  \begin{enumerate}
    \item Quando um nó quer transmitir, pode enviá-lo à taxa $R$
    \item Quando $M$ nós querem transmitir, cada um envia à taxa média $R/M$
    \item \textbf{Totalmente descentralizado:}
      \begin{itemize}
        \item nenhum nó especial para coordenar
        \item nenhuma sincronização de relógios nem compartimentos
      \end{itemize}
    \item \textbf{Simples}
  \end{enumerate}
  \begin{alertblock}{Trade-off}
    Particionamento (objetivos 1--2) vs.\ acesso aleatório (objetivos 3--4): não se podem ter todos à vez $\rightarrow$ por isso existem distintas famílias de protocolos.
  \end{alertblock}
\end{frame}

\begin{frame}{Taxonomia dos protocolos MAC}
  \begin{center}
    \begin{tikzpicture}[font=\scriptsize,
        caja/.style={draw,rounded corners=2mm,fill=aneliseAzul!13,inner sep=1mm,text width=3.6cm,align=center},
        fre/.style={thick,->},
      ]
      \node[caja](root){Protocolos MAC};
      \node[caja,left=3cm of root,yshift=0.2cm](part){Particionamento do canal\\(TDMA, FDMA, CDMA)};
      \node[caja,below=1.5cm of root](ale){Acesso aleatório\\(ALOHA, S-ALOHA, CSMA,\\CSMA/CD, CSMA/CA)};
      \node[caja,right=3cm of root,yshift=0.2cm](turn){Passagem de turno\\(Polling, Token)};
      \draw[fre] (root.south west) to[bend left=10] (part.east);
      \draw[fre] (root.south) -- (ale.north);
      \draw[fre] (root.south east) to[bend right=10] (turn.west);
    \end{tikzpicture}\\[1mm]
    {\scriptsize \textbf{Figura 1:} classificação dos protocolos MAC em particionamento, acesso aleatório e passagem de turno.}
  \end{center}
  \begin{itemize}
    \item \textbf{Particionamento:} divide o canal em pedaços (tempo, frequência, código); aloca um pedaço a cada nó
    \item \textbf{Acesso aleatório:} o canal não se divide; permitem-se colisões e ``recupera-se''
    \item \textbf{Passagem de turno:} os nós transmitem nos seus turnos; com mais volume podem usar turnos mais longos
  \end{itemize}
\end{frame}

\section{TDMA e FDMA}

\begin{frame}{TDMA --- Time Division Multiple Access}
  \begin{itemize}
    \item acesso ao canal por ``turnos''
    \item cada estação controla um \textbf{slot} (compartimento) de tamanho fixo (tempo de transmissão de um pacote) em cada turno
    \item divide o tempo em \textbf{quadros temporais} com $N$ slots ($N=$ número de computadores)
    \item \alert{slots não usados se desperdiçam}
  \end{itemize}
  \begin{center}
    \begin{tikzpicture}[font=\scriptsize,
        slot/.style={draw,minimum width=1cm,minimum height=0.7cm,anchor=south west,inner sep=0.5mm},
      ]
      \node[slot,fill=aneliseVerde!40](s0){};
      \node[slot,fill=aneliseGris!15,right=0.1cm of s0](s1){};
      \node[slot,fill=aneliseVerde!40,right=0.1cm of s1](s2){};
      \node[slot,fill=aneliseVerde!40,right=0.1cm of s2](s3){};
      \node[slot,fill=aneliseGris!15,right=0.1cm of s3](s4){};
      \node[slot,fill=aneliseGris!15,right=0.1cm of s4](s5){};
      \node[below=1mm of s0,font=\tiny]{1};
      \node[below=1mm of s1,font=\tiny]{2};
      \node[below=1mm of s2,font=\tiny]{3};
      \node[below=1mm of s3,font=\tiny]{4};
      \node[below=1mm of s4,font=\tiny]{5};
      \node[below=1mm of s5,font=\tiny]{6};
    \end{tikzpicture}\\[1mm]
    {\scriptsize \textbf{Figura 2:} divisão do tempo em slots fixos; os slots 2, 5 e 6 ficam vazios.}
  \end{center}
  \begin{exampleblock}{Exemplo}
    Rede local com 6 estações; as 1, 3 e 4 têm pacotes; os slots 2, 5 e 6 ficam vazios e se desperdiçam.
  \end{exampleblock}
\end{frame}

\begin{frame}{FDMA --- Frequency Division Multiple Access}
  \begin{itemize}
    \item o espectro do canal se divide em \textbf{bandas de frequência}
    \item cada estação recebe uma banda
    \item \alert{tempo de transmissão não usado nas bandas se desperdiça}
  \end{itemize}
  \begin{center}
    \begin{tikzpicture}[font=\scriptsize,
        banda/.style={draw,minimum width=6cm,minimum height=0.55cm,anchor=north west,inner sep=0.5mm},
      ]
      \node[banda,fill=aneliseVerde!40](f1){est. 1 (dado)};
      \node[banda,fill=aneliseGris!15,below=0.1cm of f1](f2){est. 2 (vazia)};
      \node[banda,fill=aneliseVerde!40,below=0.1cm of f2](f3){est. 3 (dado)};
      \node[banda,fill=aneliseVerde!40,below=0.1cm of f3](f4){est. 4 (dado)};
      \node[banda,fill=aneliseGris!15,below=0.1cm of f4](f5){est. 5 (vazia)};
      \node[banda,fill=aneliseGris!15,below=0.1cm of f5](f6){est. 6 (vazia)};
      \node[right=2mm of f1.east,font=\tiny]{frequência};
    \end{tikzpicture}\\[1mm]
    {\scriptsize \textbf{Figura 3:} divisão do espectro em bandas de frequência; bandas sem tráfego ficam ociosas.}
  \end{center}
\end{frame}

\section{Acesso aleatório: ALOHA}

\begin{frame}{Protocolos de acesso aleatório}
  \begin{block}{Ideias base}
    \begin{itemize}
      \item quando o nó tem um pacote: transmite a toda a taxa $R$
      \item não há regra de coordenação a priori
      \item duas ou mais transmissões simultâneas $\rightarrow$ \textbf{colisão}
    \end{itemize}
  \end{block}
  \begin{block}{O protocolo especifica}
    \begin{itemize}
      \item como detectar colisões
      \item como se recuperar (retransmissões atrasadas)
    \end{itemize}
  \end{block}
  \textbf{Exemplos:} slotted ALOHA, ALOHA puro, CSMA, CSMA/CD, CSMA/CA.
\end{frame}

\begin{frame}{Slotted ALOHA: operação}
  \textbf{Suposições:}
  \begin{itemize}
    \item todos os frames do mesmo tamanho
    \item tempo dividido em slots (tempo de transmissão de um frame)
    \item os nós transmitem só ao início de slot; estão sincronizados
    \item se 2 ou mais transmitem no slot, todos detectam a colisão
  \end{itemize}
  \textbf{Operação:}
  \begin{itemize}
    \item quando um nó obtém um novo frame, transmite no próximo slot
    \item sem colisão: pode enviar o novo frame no próximo slot
    \item com colisão: retransmite em cada slot subsequente com probabilidade $p$ até sucesso
  \end{itemize}
  \begin{center}
    \begin{tikzpicture}[font=\scriptsize,
        slot/.style={draw,minimum width=1cm,minimum height=0.65cm,anchor=south west,inner sep=0.5mm},
      ]
      \node[slot,fill=aneliseLaranja!40](s0){};
      \node[slot,fill=aneliseGris!15,right=0.1cm of s0](s1){};
      \node[slot,fill=aneliseLaranja!40,right=0.1cm of s1](s2){};
      \node[slot,fill=aneliseGris!15,right=0.1cm of s2](s3){};
      \node[slot,fill=aneliseVerde!45,right=0.1cm of s3](s4){};
      \node[slot,fill=aneliseGris!15,right=0.1cm of s4](s5){};
      \node[below=1mm of s0,font=\tiny]{C};
      \node[below=1mm of s1,font=\tiny]{E};
      \node[below=1mm of s2,font=\tiny]{C};
      \node[below=1mm of s4,font=\tiny]{S};
      \node[below=1mm of s5,font=\tiny]{E};
    \end{tikzpicture}\\[1mm]
    {\scriptsize \textbf{Figura 4:} sequência de slots com colisões (C), slot vazio (E) e sucesso (S).}
  \end{center}
  \begin{center}
    \scriptsize C = colisão, E = slot vazio, S = sucesso
  \end{center}
\end{frame}

\begin{frame}{Eficiência do Slotted ALOHA}
  \begin{itemize}
    \item eficiência $=$ fração de slots com sucesso quando há muitos nós com muitos frames
    \item $N$ nós com muitos frames; cada um transmite no slot com probabilidade $p$
    \item prob.\ de sucesso de um nó $= p(1-p)^{N-1}$
    \item prob.\ de sucesso de qualquer nó $= Np(1-p)^{N-1}$
    \item ótimo $p^*$ maximiza $Np(1-p)^{N-1}$
    \item limite quando $N\to\infty$: $\dfrac{1}{e} = 0{,}37$
  \end{itemize}
  \begin{alertblock}{Conclusão}
    No máximo: uso do canal para dados úteis: $\mathbf{37\%}$ do tempo!
  \end{alertblock}
\end{frame}

\begin{frame}{ALOHA puro (unslotted)}
  \begin{itemize}
    \item operação mais simples: \textbf{sem sincronização}
    \item o pacote é enviado sem esperar início de slot
    \item a probabilidade de colisão aumenta:
      \begin{itemize}
        \item pacote enviado em $t_0$ colide com outros enviados em $[t_0-1,\ t_0+1]$
      \end{itemize}
  \end{itemize}
  \begin{block}{Eficiência}
    \begin{align*}
      P(\text{sucesso de um nó}) &= p\,(1-p)^{N-1}\,(1-p)^{N-1} = p(1-p)^{2(N-1)}
    \end{align*}
    com o $p$ ótimo e $n\to\infty$: $\dfrac{1}{2e} = 0{,}18$
  \end{block}
  \begin{alertblock}{Conclusão}
    No máximo: $\mathbf{18\%}$ do tempo. Pior que slotted ALOHA.
  \end{alertblock}
\end{frame}

\section{CSMA, CSMA/CD}

\begin{frame}{CSMA --- Carrier Sense Multiple Access}
  \begin{block}{Ideia}
    \textbf{Escuta antes de transmitir:}
    \begin{itemize}
      \item se o canal parece vazio $\rightarrow$ transmite o pacote
      \item se o canal está ocupado $\rightarrow$ adia a transmissão
    \end{itemize}
    \emph{``Não interrompa os demais!''}
  \end{block}
  \begin{block}{Colisões no CSMA}
    \begin{itemize}
      \item o atraso de propagação implica que dois nós podem não ouvir a transmissão do outro
      \item colisão: todo o tempo de transmissão do pacote se desperdiça
      \item \textbf{importância da distância e do atraso de propagação}
    \end{itemize}
  \end{block}
\end{frame}

\begin{frame}{CSMA/CD --- Collision Detection}
  \begin{itemize}
    \item detecção de portadora e deferência como CSMA
    \item colisões detectadas em tempo mais curto
    \item transmissões com colisão são interrompidas $\rightarrow$ menos desperdício
  \end{itemize}
  \begin{block}{Como se detecta a colisão?}
    \begin{itemize}
      \item \textbf{fácil em LANs cabeada}: medição da intensidade do sinal, comparação do transmitido e recebido
      \item \textbf{difícil em LANs sem fio}: o receptor está apagado enquanto transmite!
    \end{itemize}
  \end{block}
  \begin{exampleblock}{Ethernet usa CSMA/CD}
    \begin{itemize}
      \item sem slots; \emph{carrier sense}: não transmite se detecta outra estação
      \item \emph{collision detection}: aborta ao detectar outra transmissão
      \item \emph{random access}: antes de retransmitir espera um período aleatório
    \end{itemize}
  \end{exampleblock}
\end{frame}

\begin{frame}{Algoritmo CSMA/CD da Ethernet}
  \begin{enumerate}
    \item A estação recebe um datagrama da camada de rede e cria um frame
    \item Se detecta canal livre, começa a transmitir; se ocupado, espera até que fique livre
    \item Se transmite o frame completo sem detectar outra transmissão $\rightarrow$ ``sucesso!''
    \item Se detecta outra transmissão enquanto transmite $\rightarrow$ aborta e envia \emph{jam signal} (48 bits)
    \item Após abortar: \emph{exponential backoff}: após a $m$-ésima colisão, seleciona $K\in\{0,1,\ldots,2^m-1\}$, espera $K\cdot 512$ bit-times e volta ao passo 2
  \end{enumerate}
  \begin{exampleblock}{Valores típicos}
    Bit time (10~Mbps) $= 0{,}1\,\mu$s; $K=1023 \Rightarrow$ espera $\approx 50$~ms.
    Backoff: primeira colisão $K\in\{0,1\}$; segunda $\{0,1,2,3\}$; após 10 ou mais $\{0,\ldots,1023\}$.
  \end{exampleblock}
\end{frame}

\begin{frame}{Eficiência do CSMA/CD}
  \begin{itemize}
    \item $t_{\text{prop}}$ = propagação máxima entre 2 nós na LAN
    \item $t_{\text{trans}}$ = tempo de transmissão de um frame de tamanho máximo
  \end{itemize}
  \begin{block}{Ethernet 10~Mbps}
    $t_{\text{prop}} = 51{,}2\,\mu$s, $t_{\text{trans}} = 1{,}2$~ms $\Rightarrow$ eficiência $\approx 82{,}6\%$.
  \end{block}
  \begin{itemize}
    \item tende a 1 quando $t_{\text{prop}}\to 0$ ou $t_{\text{trans}}\to\infty$
    \item \textbf{muito melhor que ALOHA}, e segue sendo descentralizado, simples e barato
  \end{itemize}
  \begin{alertblock}{Nota atualizada}
    Ethernet moderno (switches) é full-duplex; CSMA/CD fica como material histórico.
  \end{alertblock}
\end{frame}

\section{Passagem de turno}

\begin{frame}{Protocolos MAC com passagem de turno}
  \begin{columns}
    \begin{column}{0.5\textwidth}
      \textbf{Por que outra família?}
      \begin{itemize}
        \item particionamento: eficiente com carga alta, ineficiente com carga baixa (um nó só recebe $1/N$)
        \item acesso aleatório: eficiente com carga baixa (um nó usa tudo), colisões com carga alta
        \item \textbf{passagem de turno}: busca o melhor de ambos
      \end{itemize}
    \end{column}
    \begin{column}{0.5\textwidth}
      \textbf{Polling}
      \begin{itemize}
        \item nó-mestre ``convida'' os escravos a transmitir um a um
        \item problemas: overhead de polling, latência, ponto único de falha (mestre)
      \end{itemize}
      \textbf{Token passing}
      \begin{itemize}
        \item um token passa de nó a nó sequencialmente (mensagem token)
        \item problemas: overhead de token, latência, ponto único de falha (token)
      \end{itemize}
    \end{column}
  \end{columns}
  \begin{center}
    \begin{tikzpicture}[font=\scriptsize,
        n/.style={draw,circle,minimum size=6mm,fill=aneliseAzul!15,inner sep=0.4mm},
      ]
      \node[n](a){};
      \node[n,right=1.2cm of a](b){};
      \node[n,right=1.2cm of b](c){};
      \node[n,right=1.2cm of c](d){};
      \draw[-Stealth,thick,aneliseLaranja] (a) -- node[midway,above=1mm,font=\tiny]{token} (b);
      \draw[thick,aneliseGris] (b) -- (c);
      \draw[thick,aneliseGris] (c) -- (d);
      \draw[thick,aneliseGris] (d.south) to[bend right=60] (a.south);
    \end{tikzpicture}\\[1mm]
    {\scriptsize \textbf{Figura 5:} passagem de turno com o token percorrendo os nós em sequência.}
  \end{center}
\end{frame}

\section{CSMA/CA e 802.11}

\begin{frame}{Wireless Ethernet: por que não CSMA/CD?}
  \begin{columns}
    \begin{column}{0.55\textwidth}
      \textbf{802.3 (Ethernet)}: CSMA/CD
      \begin{itemize}
        \item Carrier Sense + Collision Detect
      \end{itemize}
      \textbf{802.11 (WLAN)}: CSMA/CA (Collision Avoidance)
      \begin{itemize}
        \item tanto A como C ``sentem'' o canal desocupado ao mesmo tempo e ambos enviam
        \item com CSMA/CD a colisão pode se detectar no emissor; \alert{por que não com CSMA/CA?}
      \end{itemize}
    \end{column}
    \begin{column}{0.45\textwidth}
      \begin{block}{Razões}
        \begin{itemize}
          \item comunicação \textbf{half-duplex}
          \item grande diferença entre potência de transmissão e recepção
          \item meio sem fio
        \end{itemize}
      \end{block}
      \begin{center}
        \begin{tikzpicture}[font=\scriptsize,
            n/.style={draw,circle,minimum size=6mm,fill=aneliseAzul!15,inner sep=0.4mm},
            ap/.style={draw,circle,minimum size=6mm,fill=aneliseLaranja!30,inner sep=0.4mm},
          ]
          \node[n](A){A};
          \node[ap,right=1.4cm of A](B){B};
          \node[n,right=1.4cm of B](C){C};
          \draw[thick,aneliseLaranja] (A) -- (B);
          \draw[thick,aneliseLaranja] (C) -- (B);
        \end{tikzpicture}
      \end{center}
    \end{column}
  \end{columns}
\end{frame}

\begin{frame}{IEEE 802.11 MAC: CSMA/CA}
  \begin{block}{Emissor 802.11 CSMA}
    \begin{itemize}
      \item se o canal é sentido vazio por \textbf{DIFS} segundos $\rightarrow$ envia o frame completo (sem detecção de colisão)
      \item se o canal é sentido ocupado $\rightarrow$ \textbf{binary backoff}
    \end{itemize}
  \end{block}
  \begin{block}{Receptor 802.11 CSMA}
    \begin{itemize}
      \item se o frame é recebido OK $\rightarrow$ retorna \textbf{ACK} após \textbf{SIFS} segundos
    \end{itemize}
  \end{block}
  \begin{exampleblock}{NAV --- Network Allocation Vector}
    \begin{itemize}
      \item o frame 802.11 tem um campo com o tempo de transmissão
      \item outros nós (escutando) diferem o acesso por \textbf{NAV} unidades de tempo
    \end{itemize}
  \end{exampleblock}
\end{frame}

\begin{frame}{Problema do terminal oculto}
  \begin{columns}
    \begin{column}{0.5\textwidth}
      \begin{itemize}
        \item grande diferença na força do sinal, problemas do canal físico (\emph{shadowing})
        \item duas estações não podem se comunicar na mesma área
      \end{itemize}
      \textbf{Terminal oculto:}
      \begin{enumerate}
        \item A e C não se ouvem
        \item A transmite a B
        \item C quer transmitir a B; C não pode receber A; C sente o meio desocupado (\emph{carrier sense fails})
        \item colisão em B
        \item A não pode detectar a colisão (\emph{collision detection fails})
        \item A está ``oculto'' para C
      \end{enumerate}
    \end{column}
    \begin{column}{0.5\textwidth}
      \begin{center}
        \begin{tikzpicture}[font=\scriptsize,
            n/.style={draw,circle,minimum size=7mm,fill=aneliseAzul!15,inner sep=0.5mm},
          ]
          \node[n](A){A};
          \node[n,right=1.8cm of A](B){B};
          \node[n,right=1.8cm of B](C){C};
          \draw[->,thick,aneliseLaranja] (A) -- node[midway,above=1mm,font=\tiny]{TX} (B);
          \draw[->,thick,aneliseLaranja] (C) -- node[midway,below=1mm,font=\tiny]{TX} (B);
          \draw[dashed,thick,aneliseGris] (A) to[bend left=25] node[midway,below=4mm,font=\tiny]{não se ouvem} (C);
          \node[below=2mm of B,font=\tiny]{colisão em B};
        \end{tikzpicture}\\[1mm]
        {\scriptsize \textbf{Figura 6:} problema do terminal oculto: A e C não se ouvem e colidem em B.}
      \end{center}
      \begin{alertblock}{Consequência}
        802.11 não usa detecção de portadora pura; colisões só se deduzem após a transmissão (CRC / ausência de ACK).
      \end{alertblock}
    \end{column}
  \end{columns}
  \note{Sugestão de demonstração em sala: escolha 3 alunos como A, B e C; B fica no centro e A e C nas extremidades, sem se enxergarem. Peça que A e C transmitam ao mesmo tempo para B e observe que a colisão em B não é detectada nem por A nem por C.}
\end{frame}

\begin{frame}{Acesso ao canal: IFS e backoff}
  \begin{columns}
    \begin{column}{0.5\textwidth}
      \textbf{DIFS --- Distributed Inter Frame Space}
      \begin{itemize}
        \item uma estação pode transmitir só se sente o canal desocupado por um tempo \textbf{DIFS}
        \item objetivo: separar dados e ACK com IFS distintos (SIFS $<$ DIFS)
      \end{itemize}
      \textbf{Backoff}
      \begin{itemize}
        \item quando o canal inicialmente está ocupado, a estação deve diferir; quando fica desocupado por DIFS, difere ainda por um tempo aleatório (\emph{backoff time})
        \item \textbf{slotted backoff}: extrair $K\in(0,W-1)$; decrementar cada slot-time
        \item \textbf{backoff freezing}: o contador ``se congela'' enquanto o canal estiver ocupado
      \end{itemize}
    \end{column}
    \begin{column}{0.5\textwidth}
      \textbf{Regras de backoff}
      \begin{itemize}
        \item primeiro: $K\in(0,\mathrm{CW}_{\min})$
        \item sem sucesso: $K\in(0, 2(\mathrm{CW}_{\min}+1)-1)$
        \item novamente: $K\in(0, 2^2(\mathrm{CW}_{\min}+1)-1)$
        \item após $m$ falhas: $K\in(0, 2^m(\mathrm{CW}_{\min}+1)-1)$
      \end{itemize}
      \begin{exampleblock}{802.11 clássico}
        $\mathrm{CW}_{\min}=31$, $\mathrm{CW}_{\max}=1023$ ($m=5$). Resultado: \textbf{backoff exponencial}.
      \end{exampleblock}
      \begin{center}
        \begin{tikzpicture}[font=\scriptsize]
          \node[draw,fill=aneliseVerde!12,rounded corners=2mm,text width=5.6cm,align=center](res){Referência: G.\ Bianchi, ``Performance Analysis of the IEEE 802.11 DCF'', IEEE JSAC 18(3), 2000.};
        \end{tikzpicture}
      \end{center}
    \end{column}
  \end{columns}
\end{frame}

\section{RTS/CTS}

\begin{frame}{RTS/CTS --- Request-To-Send / Clear-To-Send}
  \begin{itemize}
    \item \textbf{4-way handshaking} no lugar do 2-way básico
    \item introduzido por duas razões:
      \begin{enumerate}
        \item resolve o problema do \textbf{terminal oculto}
        \item aumenta o throughput com pacotes longos
      \end{enumerate}
  \end{itemize}
  \begin{block}{Mecanismo}
    \begin{itemize}
      \item o transmissor envia \textbf{RTS} curto
      \item o receptor responde \textbf{CTS}
      \item CTS reserva o canal para o transmissor, notificando as outras estações (possivelmente ocultas)
      \item RTS e CTS carregam o tempo que o canal estará ocupado $\rightarrow$ as outras atualizam o \textbf{NAV} (\emph{virtual carrier sensing})
    \end{itemize}
  \end{block}
  \begin{exampleblock}{Prós/contras}
    RTS/CTS curtos $\rightarrow$ colisões menos prováveis e de menor duração. Overhead maior.
  \end{exampleblock}
\end{frame}

\begin{frame}{Collision Avoidance: troca RTS--CTS}
  \begin{center}
    \begin{tikzpicture}[font=\scriptsize,
        n/.style={draw,circle,minimum size=7mm,fill=aneliseAzul!15,inner sep=0.5mm},
      ]
      \node[n](A){A (origem)};
      \node[n,right=2.2cm of A](B){B (destino)};
      \node[n,above=1.5cm of B](C){C (outros)};
      \node[n,below=1.5cm of B](D){D (outros)};
      \draw[->,thick,aneliseLaranja] (A) -- node[midway,above=1mm,font=\tiny]{RTS} (B);
      \draw[->,thick,aneliseVerde] (B) -- node[midway,below=1mm,font=\tiny]{CTS} (A);
      \draw[dashed,->,thick,aneliseGris] (B.north) to[bend left=15] (C.south);
      \draw[dashed,->,thick,aneliseGris] (B.south) to[bend right=15] (D.north);
      \node[draw,fill=aneliseVerde!12,rounded corners=2mm,below=6mm of A,text width=4.6cm,align=center](nav){C e D atualizam NAV e diferem};
    \end{tikzpicture}\\[1mm]
    {\scriptsize \textbf{Figura 7:} troca RTS--CTS: o CTS de B é escutado por C e D, que atualizam o NAV e adiam o acesso.}
  \end{center}
  \begin{itemize}
    \item evita colisões com estações ocultas
    \item NAV $=$ acesso diferido (\emph{virtual carrier sensing})
    \item 802.11 permite CSMA, CSMA/CA com reservas, e PCF (polling desde o AP)
  \end{itemize}
  \note{Relembre o problema do terminal oculto: o RTS de A pode não chegar a C, mas o CTS de B é ouvido por toda a vizinhança, incluindo C e D, que então atualizam o NAV.}
\end{frame}

\begin{frame}{Throughput do RTS/CTS e relação IFS}
  \begin{columns}
    \begin{column}{0.5\textwidth}
      \textbf{Quando convém RTS/CTS?}
      \begin{itemize}
        \item RTS/CTS contras: overhead maior
        \item RTS/CTS prós: duração do tempo de colisão reduzida
        \item convém para \textbf{pacotes longos} e \textbf{muitos terminais} (colisão!)
        \item mais robusto diante de maior número de usuários e configuração de $\mathrm{CW}_{\min}$
      \end{itemize}
    \end{column}
    \begin{column}{0.5\textwidth}
      \textbf{Relação entre IFS}
      \begin{center}
        \begin{tikzpicture}[font=\scriptsize]
          \node[draw,fill=aneliseGris!14,rounded corners=2mm,text width=4.4cm,align=center](ifs){SIFS $<$ PIFS $<$ DIFS};
          \node[draw,fill=aneliseVerde!12,rounded corners=2mm,below=4mm of ifs,text width=4.6cm,align=center](pcf){PIFS usado pelo PCF\\ (point coordination function)};
          \node[draw,fill=aneliseLaranja!12,rounded corners=2mm,below=4mm of pcf,text width=4.8cm,align=center](nota){PCF nunca foi implementado};
        \end{tikzpicture}
      \end{center}
      \begin{itemize}
        \item PCF: esquema baseado em polling desde o AP (opcional, quase nunca implementado)
      \end{itemize}
    \end{column}
  \end{columns}
\end{frame}

\begin{frame}{Por que utilizar NAV?}
  \begin{center}
    \begin{tikzpicture}[font=\scriptsize,
        n/.style={draw,circle,minimum size=7mm,fill=aneliseAzul!15,inner sep=0.5mm},
      ]
      \node[n](TX){TX};
      \node[n,right=2.2cm of TX](RX){RX};
      \node[n,above=1.6cm of TX](C){C};
      \draw[->,thick,aneliseLaranja] (TX) -- node[midway,above=1mm,font=\tiny]{frame} (RX);
      \draw[->,thick,aneliseLaranja] (TX) -- node[midway,left=1mm,font=\tiny]{frame} (C);
      \draw[<->,thick,aneliseVerde] (RX) -- node[midway,above=2mm,font=\tiny]{ACK} (C);
      \node[draw,fill=aneliseGris!12,rounded corners=2mm,below=5mm of RX,text width=5.2cm,align=center](ex){C recebe o frame de TX; C não alcança RX; mas C seta NAV e protege o ACK de RX};
    \end{tikzpicture}\\[1mm]
    {\scriptsize \textbf{Figura 8:} o NAV permite que C proteja o ACK de RX mesmo sem escutar RX (carrier sensing virtual).}
  \end{center}
  \note{Compare o \emph{carrier sensing} físico (escutar o canal) com o virtual (NAV): o NAV é definido pelo campo Duration do cabeçalho 802.11 e protege o ACK mesmo quando a estação não alcança o receptor.}
\end{frame}

\section{PCF e resumo}

\begin{frame}{PCF --- Point Coordination Function}
  \begin{columns}
    \begin{column}{0.55\textwidth}
      \begin{itemize}
        \item mecanismo de acesso baseado em \textbf{turno} (\emph{polling})
        \item o acesso decide o \textbf{point coordinator} (PC): tipicamente o AP
        \item acesso livre de contenção (sem colisões no canal)
        \item parte \textbf{opcional} da especificação $\rightarrow$ quase nunca implementado; possível uso com 802.11e
        \item os APs enviam \textbf{beacons} em intervalos regulares (0,1~s)
      \end{itemize}
    \end{column}
    \begin{column}{0.45\textwidth}
      \textbf{Dois períodos entre beacons:}
      \begin{itemize}
        \item \textbf{CFP} (Contention Free Period): o AP envia \emph{CF-Poll} a cada estação
        \item \textbf{CP} (Contention Period): usa-se DCF
      \end{itemize}
      \begin{alertblock}{Limitações}
        \begin{itemize}
          \item transição CP$\to$CFP não determinista (espera o fim da transmissão)
          \item round robin interroga todas as estações, mesmo sem dados $\rightarrow$ problemas de escalabilidade
        \end{itemize}
      \end{alertblock}
    \end{column}
  \end{columns}
\end{frame}

\begin{frame}{Resumo dos protocolos MAC}
  \begin{block}{Como se gerencia um canal compartilhado?}
    \begin{enumerate}
      \item \textbf{Particionamento do canal} (tempo, frequência, código): TDMA, FDMA, CDMA
      \item \textbf{Particionamento aleatório (dinâmico)}: ALOHA, S-ALOHA, CSMA, CSMA/CD
        \begin{itemize}
          \item detecção de portadora: fácil em cabos, difícil em sem fio
          \item CSMA/CD: Ethernet \quad CSMA/CA: 802.11
        \end{itemize}
      \item \textbf{Passagem de turno}: polling centralizado, token
    \end{enumerate}
  \end{block}
  \begin{alertblock}{Evolução (atualizado)}
    802.11e/ax (OFDMA, QoS), 802.11s mesh, 5G NR (scheduling centralizado) são evoluções dessas ideias.
  \end{alertblock}
\end{frame}

\begin{frame}{Exercício resolvido: eficiência do Slotted ALOHA}
  \textbf{Enunciado.} Uma rede slotted ALOHA tem $N=20$ estações. Cada estação transmite em um slot com probabilidade $p=0{,}05$. Qual é a eficiência $S$ (fração de slots com sucesso)?
  \begin{block}{Solução}
    A eficiência do Slotted ALOHA é dada por
    \[
      S = N\,p\,(1-p)^{N-1}.
    \]
    Substituindo $N=20$ e $p=0{,}05$:
    \[
      S = 20 \cdot 0{,}05 \cdot (1-0{,}05)^{20-1} = 1 \cdot (0{,}95)^{19}.
    \]
    Como $(0{,}95)^{19} \approx 0{,}377$:
    \[
      \boxed{S \approx 0{,}377 \approx 37{,}7\%}
    \]
  \end{block}
  \begin{alertblock}{Comparação com o máximo teórico}
    O limite quando $N\to\infty$ (com $p^*=1/N$) é $1/e \approx 0{,}37$. Com $N=20$ e $p=0{,}05$, a eficiência $\approx 0{,}377$ já atinge praticamente o máximo teórico do protocolo.
  \end{alertblock}
  \begin{exampleblock}{Extensão: ALOHA puro}
    Sem slots, o throughput máximo é $S_{\max} = \dfrac{1}{2e} \approx 0{,}18$ --- cerca de metade do Slotted ALOHA.
  \end{exampleblock}
\end{frame}

\begin{frame}{Resumo da aula}
  \begin{itemize}
    \item \textbf{Enlaces broadcast} precisam de um protocolo de acesso múltiplo para arbitrar o canal compartilhado.
    \item \textbf{Particionamento} (TDMA, FDMA): divisão justa do canal, mas desperdício de capacidade com carga variável.
    \item \textbf{ALOHA}: acesso aleatório simples; Slotted ALOHA atinge $1/e \approx 0{,}37$ e ALOHA puro $1/(2e) \approx 0{,}18$.
    \item \textbf{CSMA}: escuta antes de transmitir; \textbf{CSMA/CD} funciona em cabos, mas não em rádio (half-duplex).
    \item \textbf{CSMA/CA (802.11)}: DIFS/SIFS, backoff exponencial, ACK e \textbf{NAV} (\emph{virtual carrier sensing}).
    \item \textbf{RTS/CTS}: resolve o problema do terminal oculto; \textbf{PCF} (polling) é opcional e quase nunca implementado.
  \end{itemize}
\end{frame}

\begin{frame}{Autoavaliação}
  \begin{enumerate}
    \item Qual é a eficiência máxima teórica do Slotted ALOHA quando o número de estações tende ao infinito?
      \begin{itemize}
        \item (a) $1/(2e) \approx 0{,}18$
        \item (b) $1/e \approx 0{,}37$
        \item (c) $1{,}0$
      \end{itemize}
      \pause
      \textbf{Gabarito:} (b)
    \item Por que o CSMA/CD não é empregado em redes sem fio 802.11?
      \begin{itemize}
        \item (a) porque o canal é compartilhado apenas por duas estações
        \item (b) porque o transmissor não consegue detectar colisões enquanto transmite (half-duplex)
        \item (c) porque o backoff exponencial não existe no rádio
      \end{itemize}
      \pause
      \textbf{Gabarito:} (b)
    \item Qual é a função do NAV (\emph{Network Allocation Vector}) no 802.11?
      \begin{itemize}
        \item (a) definir o tamanho do pacote a ser transmitido
        \item (b) reservar o canal por um tempo, mesmo para estações que não escutam o transmissor
        \item (c) sincronizar os relógios das estações com o AP
      \end{itemize}
      \pause
      \textbf{Gabarito:} (b)
    \item O handshake RTS/CTS foi introduzido no 802.11 principalmente para:
      \begin{itemize}
        \item (a) resolver o problema do terminal oculto
        \item (b) eliminar totalmente qualquer colisão
        \item (c) aumentar a potência de transmissão
      \end{itemize}
      \pause
      \textbf{Gabarito:} (a)
  \end{enumerate}
\end{frame}

\begin{frame}{Laboratório sugerido: simulação e análise de tráfego}
  \textbf{Opção 1 --- Simulação no ns-3}
  \begin{itemize}
    \item Usar o módulo \texttt{wifi} (ex.: \texttt{WifiHelper} e \texttt{YansWifiChannelHelper}) com o modelo 802.11;
    \item Cenário: $N$ estações enviando tráfego UDP para o AP;
    \item Medir \textbf{throughput} agregado e o número de \textbf{colisões} com RTS/CTS habilitado ($\mathrm{RtsCtsThreshold}=0$) e desabilitado (valor alto);
    \item Comparar os resultados para pacotes curtos e longos.
  \end{itemize}
  \textbf{Opção 2 --- Wireshark em uma rede Wi-Fi doméstica}
  \begin{itemize}
    \item Capturar frames 802.11 (modo monitor ou na rede do AP);
    \item Observar \textbf{DIFS}, \textbf{backoff}, a troca de \textbf{ACK} e o campo \textbf{Duration} (NAV) do cabeçalho;
    \item Identificar retransmissões e o efeito do \emph{carrier sense}.
  \end{itemize}
\end{frame}

\section{Bibliografia}

\begin{frame}{Bibliografia}
  \begin{itemize}
    \item \textbf{\kurose{}} --- Capítulo 5 (seção 5.3)
  \end{itemize}
  \begin{block}{Adicional}
    \begin{itemize}
      \item G.\ Bianchi, ``Performance Analysis of the IEEE 802.11 DCF'', \emph{IEEE JSAC} 18(3):535--547, 2000
      \item Jurdak, Lopes, Baldi, ``A Survey, Classification, and Comparative Analysis of MAC Protocols for Ad Hoc Networks'', \emph{IEEE Comm.\ Surveys \& Tutorials} 6(1):2--16, 2004
      \item Heusse, Rousseau, Berger-Sabbatel, Duda, ``Performance anomaly of 802.11b'', \emph{IEEE INFOCOM} 2003
    \end{itemize}
  \end{block}
\end{frame}

\begin{frame}{Leitura recomendada}
  \begin{itemize}
    \item \textbf{\kurose{}} --- Capítulo 5 (seção 5.3)
    \item \textbf{\stallings{}} --- Capítulos sobre MAC sem fio e padrão 802.11
    \item \textbf{\forouzan{}} --- Capítulo 12 (``Multiple Access'')
    \item G.\ Bianchi, ``Performance Analysis of the IEEE 802.11 DCF'', \emph{IEEE JSAC} 18(3):535--547, 2000
  \end{itemize}
  \begin{block}{Sugestão de leitura}
    Aprofunde a análise de desempenho do CSMA/CA em Bianchi (2000) e o comportamento sob assimetria de taxas em Heusse et al.\ (2003).
  \end{block}
\end{frame}

\end{document}